\documentclass[11pt,a4paper]{article}
\usepackage[margin=24mm]{geometry}
\usepackage{amsmath,amssymb,amsthm,booktabs,microtype}
\usepackage[T1]{fontenc}
\usepackage{lmodern}
\usepackage[colorlinks=true,linkcolor=blue,urlcolor=blue,citecolor=blue]{hyperref}
\hypersetup{pdftitle={Subset Sums: a unified loss and hole proof from 98 to 97},pdfauthor={Sum-of-Subset research project}}
\newcommand{\FS}{\operatorname{FS}}
\newcommand{\Z}{\mathbb Z}
\newcommand{\T}{\mathcal T}
\newtheorem{theorem}{Theorem}[section]
\newtheorem{lemma}[theorem]{Lemma}
\newtheorem{proposition}[theorem]{Proposition}
\title{Subset Sums in a Fourfold Sumset\\\large A unified loss and hole proof from 98 to 97}
\author{Sum-of-Subset research project}
\date{16 October 2026}
\begin{document}
\maketitle
\begin{abstract}
We exclude cardinality 98, continuing the independently audited bound to
$N\le97$. The proof replaces endpoint-specific full-support rigidity by
one span budget valid with holes, retains a single shared hole parameter
in all model rows, and classifies exceptional translates by their loss
rather than their outside endpoint. A refined lattice comparison closes
a gap in directly reusing the previous coarse criterion. The complete
finite reduction consists of 45,928 models, covered by 3,750 symbolic
quadratic interval checks using 209 newly generated integer certificates.
Every integer model was also checked separately against the symbolic
formulas. The independently closed range is $71\le N\le97$; the stronger
archived claim $N\le84$ remains outside the complete independent audit.
\end{abstract}
\tableofcontents

\section{Statement and proof organization}
For a finite real set $A$, let
\[
 \FS(A)=\left\{\sum_{a\in X}a:X\subseteq A\right\},\qquad
 4A=A+A+A+A.
\]
The empty sum is included and repeated summands in $4A$ are allowed.
Let $N$ denote the maximum feasible cardinality.
\begin{theorem}\label{thm:main}
If $\FS(A)\subseteq4A$, then $|A|\le97$.
\end{theorem}
The existing feasible 71-element set is unchanged, so $71\le N\le97$.
Our foundation is the independently checked chain
$115\to103\to100\to99\to98$ \cite{base103,base100,base99,base98}.
We use its integralization, covering, peeling, inverse, nonprogression,
and effective-hole lemmas. The archived upper-86 report \cite{archive}
contains the strategy for the present endpoint. We reconstruct the
necessary reduction and generate fresh certificates; no archived endpoint
certificate file is read. This is an advance of the audited chain, not a
numerical improvement on the archived upper bound 84.

The optimization is mathematical as well as computational:
\begin{enumerate}
\item The universal span estimate already proved in \cite{base103} is
combined with the budget before losses and holes are separated. No
zero-defect classification or full-support reduction is needed here.
\item A single variable $1\le z\le14$ is shared by the loss, span, mass,
and sumset rows. The certificate proves inconsistency for every real
value in this interval, so it covers every possible integer hole count.
\item Core quantiles apply whenever a retained translate has more than
four hull points. With one empty unit, at most one translate can violate
this threshold. Designating that translate replaces enumeration of its
outside endpoint or of multiple outside centers.
\item The nested envelopes are imposed uniformly, including the full
hull endpoint. Valid slot-count caps and the coarse double-sumset bound
are retained. The previous exact quadratic interval engine verifies all
integer side sizes without repeating the same matrix combination.
\end{enumerate}

\section{A refined lattice comparison at size 98}
A hypothetical feasible 98-element set contains zero: otherwise adjoining
zero produces a forbidden feasible 99-element set. Apply the inherited
integralization and remove the outer four elements on each side to obtain
$D$, then peel twice to $E=P\cup\{0\}\cup(-Q)$.
Writing $T_r=r(r+1)/2$, the parameters are
\begin{gather*}
 d=90,\quad e=82,\quad m=a+b=81,\quad a=|P|,\quad b=|Q|,\\
 \{a,b\}=\{40-k,41+k\},\quad I=k(k+1)=(a-40)(a-41),\\
 \Lambda=1+T_a+T_b=1682+I,\qquad\Omega=395-I.
\end{gather*}
The one-sided alternative is excluded by $1+T_{82}>24\cdot98+77$.
The paired bootstrap $q^2\le96e+1073-e^2$, $q\equiv e-1\pmod2$,
gives at least 17 opposite pairs, hence $a,b\ge17$.

For a positive set $X$ of size $r$ define
\[
 \mu(X)=|\FS(X)|-1-T_r,\qquad
 \delta_0(X)=|\{x\in X:2x\notin R(X\cup\{0\})\}|,
 \quad\rho(X)=\mu(X)+2-2\delta_0(X),
\]
where $R$ is the restricted two-sum set. Put
$\rho=\rho(P)+\rho(Q)$, $\mu=\mu(P)+\mu(Q)$ and $s=|\FS(E)|$.
The inherited nonnegative slacks satisfy
\begin{gather}
 \kappa=s-\Lambda-\mu,\qquad
 c=2|2D|+18d+85-|\FS(D)|,\notag\\
 \rho+2\theta+c+\T=\Omega,\qquad
 \T=\kappa+\ell_1+\ell_2+g.\label{eq:budget}
\end{gather}
All inverse gates $T_{\min(a,b)}-2-\Omega>0$ hold; their minimum over
the 40 possible oriented pairs is 163. The simultaneous inverse theorem
therefore puts $P,Q$ on a common integer lattice and gives
\begin{equation}\label{eq:coreholes}
 h_P,h_Q\le\kappa,\quad z=h_P+h_Q-\kappa,
 \quad0\le z\le\min(h_P,h_Q)\le\kappa,
 \quad h_X<|\FS(X)|-1.
\end{equation}
Here $h_X$ is the number of missing integers in the support hull of $X$,
and $z$ is the corresponding count for $E$.

\begin{lemma}[Universal span estimate]\label{lem:span}
If $X\subset\Z_{>0}$ has $r\ge15$ elements, then
$\max X\le r+\rho(X)+h_X+2$.
\end{lemma}
\begin{proof}
Delete its largest element $r+q$, use primes for the preceding prefix,
and let $\nu$ count old diagonal defects destroyed at the last step.
The new maximum is defective and
$\mu-\mu'=q-h_X+h_{X'}$, so exactly
\[
 \rho(X)=\rho(X')+q-h_X+h_{X'}-2+2\nu.
\]
The prefix has at least 14 elements, hence $\rho(X')\ge0$; also
$h_{X'},\nu\ge0$. Rearrangement proves the inequality.
Neither primitivity nor full subset-sum support is required.
\end{proof}

Set $R_p=\max P$, $R_q=\max Q$. The lemma gives
$R_p+R_q\le m+\rho+\kappa+z+4$.
Before asserting integrality of the peeled layers, use
$|R(E)|\le2(R_p+R_q)+1$ and the at most $4e+6$ new restricted pairs to get
\begin{equation}\label{eq:inner}
 |\FS(D_1)|-2s
 \le C+2\rho-\mu+\kappa+2z+\ell_2,
 \qquad C=2m+4e+14-\Lambda=-1178-I.
\end{equation}
The old coarse bound $C+3\Omega=7-4I$ fails to be negative at balance.
We retain the terms that it discarded:
\[
 |\FS(D_1)|-2s\le
 \begin{cases}
 C+2\Omega=-388-3I,&z=0,\\
 C+3\Omega-\rho-\mu\le-121-4I,&z>0.
 \end{cases}
\]
For the first line use $\rho+\kappa+\ell_2\le\Omega$,
$\rho\le\Omega$, and $\mu\ge0$ in \eqref{eq:inner}.
For the second, $z\le\kappa$ and the same budget yield
$2\rho+\kappa+2z+\ell_2\le3\Omega-\rho$.
Positive $z$ makes both side supports nonprogressions: a progression of
step at least two has at least its cardinality minus one holes, contrary
to \eqref{eq:coreholes}. The inherited penalties give
$\rho\ge m-28=53$, $\mu\ge m-6=75$.
Both comparisons are strictly negative. An off-lattice peeled point would
produce two disjoint $s$-point translates, so $D_1$ is on the lattice.
The outer comparison has margin
\[
 \Lambda-(4d-10+\Omega)=937+2I>0,
\]
which similarly puts $D$ on the lattice. The eight outer elements are
not yet asserted integral.

The inherited overlap propagation and four-extreme defect estimate now
apply without change. Writing $w$ for the support holes of $D$ and
$G=2d-12=168$, they give
\begin{equation}\label{eq:propagation}
 w\le4z,\quad\T\ge\kappa+168-3z,
 \qquad\T\ge\max\{56,z,168-2z\}.
\end{equation}
This ordering matters: the estimates on $D$ are used only after its
integrality has been established by the refined comparison.

\section{One span budget and one hole parameter}
Write $M_p=\max D$, $M_q=-\min D$, and let $S_p,S_q$ be the positive
and negative-magnitude masses of $D$. Let $j$ be the number of empty
covering units and put $\delta=c-jd$. Combining Lemma~\ref{lem:span},
\eqref{eq:budget}, and the stronger part of \eqref{eq:propagation} gives
\begin{equation}\label{eq:joint}
 \boxed{R_p+R_q+\delta\le m+4+\Omega-G-jd+4z.}
\end{equation}
Indeed $\rho+\kappa+\delta\le\Omega-G-jd+3z$, while
$R_p+R_q\le m+4+\rho+\kappa+z$.
This single estimate replaces separate zero-core and small-defect
full-support reductions in the preceding endpoint proof.

For $z=0$, $c\le227-I$. For $1\le z\le14$,
$c\le174-I+2z\le202-I$. Thus at most two units are empty.
Set $B_0=227-90j-I$ and
$\varepsilon_p=\sum P-T_a=\mu(P)+h_P$, with the reflected definition.
Zero-hole models use
\begin{equation}\label{eq:Z}
 \delta\le B_0,\qquad R_p+R_q+\delta\le85+B_0,
 \qquad\varepsilon_p,\varepsilon_q\ge0.
\end{equation}
Positive low-hole models use one common variable $z$:
\begin{gather}
 1\le z\le14,\quad\delta-2z\le B_0-53,
 \quad R_p+R_q+\delta-4z\le85+B_0,\notag\\
 \varepsilon_p\ge a-3+z,\qquad
 \varepsilon_q\ge b-3+z.\label{eq:P}
\end{gather}
The mass and cardinality rows below use this same variable. Treating it
as real enlarges the actual integer parameter set and is safe.

\subsection{A loss threshold instead of outside-endpoint enumeration}
For a retained positive center $v$, let $l_v,u_v$ count the points of
$v+D$ below 1 and above $S_p$. Let $\lambda_v=l_v+u_v$.
Reflection defines the negative-center losses using $-D$ and $S_q$.
The covering inequality gives
\begin{equation}\label{eq:loss}
 \sum_{v\text{ retained}}\lambda_v\le\delta.
\end{equation}
Every retained center contributes; it is integral since both $D$ and its
target support are integral. Therefore $0\le\lambda_v\le89$.

The necessary positive-center endpoint rows are
\begin{center}\small
\begin{tabular}{lll}
\toprule
Type&Lower row&Upper row\\
\midrule
Ordinary&$-v-l_v\le-b-5$&$v-u_v-S_p\le-a-4$\\
Global&$-v-l_v-M_p\le-90$&$v-u_v-S_p-M_q\le-89$\\
Core&$-v-l_v-R_p\le-86$&$v-u_v-S_p-R_q\le-85$\\
\bottomrule
\end{tabular}
\end{center}
Ordinary rows hold at the relevant in-hull endpoint. Global rows hold
for every nonempty translate: its $90-l_v$ points above zero fit in
an integer interval with at most $M_p+v$ positions, and its $90-u_v$
points below the upper endpoint fit in at most $S_p-v+M_q+1$ positions.
Core rows hold if the hull count exceeds four. At an outside endpoint
such a translate must use a core point because only four points on that
side lie outside the core. Count the missing positions in the core hull,
or use the corresponding ordinary row with $R_p\ge a$, $R_q\ge b$.
This gives the displayed core constants $m+5=86$ and $m+4=85$.

With one empty unit, $\delta\le137$ at zero holes and $\delta\le112$
at positive low holes. Both are below $2(90-4)=172$.
Thus at most one retained center can have $\lambda_v\ge86$.
If none does, every center uses core rows and the valid count cap
$\lambda_v\le85$. Otherwise designate the unique large-loss center,
use global rows and $86\le\lambda_v\le89$ for it, and core rows with
$\lambda_v\le85$ for the others.
With two empty units, $\delta\le47$ or $22$, so every hull count exceeds
four. At positive low holes, $\delta<\min(a,b)+5$, the minimum outside
loss; ordinary rows therefore hold for all retained centers.
With no empty unit we simply use global rows, avoiding another
all-active case split at low holes.

\section{High holes and the flexible defect inverse}
\begin{lemma}\label{lem:flex}
Let $X$ be a positive real set of size $r\ge15$, and $2\le\beta\le10$
be an integer. If $0\le\rho(X)<(\beta-1)(r-14)$, then
\[
 \mu(X)\le\rho(X)+4\beta-2+
                   2\left\lfloor\frac{\rho(X)}{\beta-1}\right\rfloor.
\]
If $\mu(X)\ge u$ is an integer lower bound, write
$v=\max(0,u-4\beta+2)=(\beta+1)q+s$, $0\le s\le\beta$. Then
\begin{equation}\label{eq:flex}
 \rho(X)\ge L_\beta(r,u):=
 \min\{(\beta-1)(r-14),(\beta-1)q+\min(s,\beta-1)\}.
\end{equation}
\end{lemma}
\begin{proof}
Use the inherited recurrence
$\rho(X_i)\ge\rho(X_{i-1})+y_i-2$.
If all $y_i>\beta$ for $15\le i\le r$, the stated hypothesis fails.
Choose the last $i\ge15$ with $y_i\le\beta$. Since $\beta\le i-5$,
the small-increment inverse lemma makes the prefix, up to scaling,
an interval with at most $\beta$ deleted positions.
The deleted-interval defect count proved in \cite{base98} needs only
$i>\beta$ and gives at most $2\beta$ diagonal defects; the zero-deletion
case has one. Each later step adds at most one defect and raises $\rho$
by at least $\beta-1$. Hence
$r-i\le\lfloor\rho(X)/(\beta-1)\rfloor$. Substitute in
$\mu=\rho-2+2\delta_0$.
For the inverse, the least integer $t\ge0$ with
$t+2\lfloor t/(\beta-1)\rfloor\ge v$ is
$(\beta-1)q+\min(s,\beta-1)$: a full block of $\beta-1$ in $t$
adds $\beta+1$ to the expression. Separate the alternative
$\rho(X)\ge(\beta-1)(r-14)$ to obtain \eqref{eq:flex}.
\end{proof}

We use the same effective-hole witness search as \cite{base98}:
$20\le r\le64$, $7\le\gamma\le20$, $0\le C\le600$.
Every chosen $(r,C,\gamma)$ is checked by the exact integer tests
\begin{gather*}
 K_0=\max\{k\ge3:k(k-3)/2\le C+1\},\quad D_*=C+K_0,\quad E_*=D_*+1,\\
 M=\max\left\{2\gamma,1+\left\lfloor
 \frac{1+\sqrt{9+8(E_*+6\gamma)}}2\right\rfloor\right\},\quad
 J=\max\{K_0+1,2\gamma+5,M\},\\
 T_M-6\gamma>M+1+E_*,\quad
 \gamma(r-J+1)>D_*,\quad T_{J-1}-1>D_*.
\end{gather*}
A valid row gives $\mu(X)\ge r+C+1$ if
$2\gamma<h_X<|\FS(X)|-1$. For each $z\ge15$ choose the largest
available $C$ with $2\gamma<z$, breaking ties by the smaller $\gamma$.
Put
\[
 u_r=r+C+1,\qquad
 v_r=\max\{r-14,u_r+2-2r,L_2(r,u_r),\ldots,L_{10}(r,u_r)\}.
\]
Then
\begin{equation}\label{eq:highbudget}
 c\le C_k(z):=395-I-v_a-v_b-\max\{56,z,168-2z\}.
\end{equation}
The loop $15\le z\le\Omega$ is exhaustive because
$z\le\kappa\le\T\le\Omega$. Impossible negative budgets are discarded.

For all-active configurations the inherited mass-bonus alternatives,
in both orientations, are
\begin{align*}
 F&\ge F_A(a,b)+\tfrac12\varepsilon_p+\varepsilon_q,
 &F_A&=\tfrac12T_a+T_b+6a+8b+85,\\
 \text{or }F&\ge F_B(a,b)+\tfrac23\varepsilon_p+\varepsilon_q,
 &F_B&=\tfrac23T_a+T_b+6a+(22b+236)/3,
\end{align*}
where $F=S_p+S_q+c-2(M_p+M_q)-2(R_p+R_q)\le1746+4z$.
Use $\varepsilon_p\ge u_a+z$, $\varepsilon_q\ge u_b+z$.
There are 1,574 surviving oriented high-hole states. Their minimum active
margin is 35 and their maximum $C_k(z)$ is 91. Thus every all-active
high-hole state is excluded, and at most one unit is empty.
Only four oriented states permit an empty unit:
\begin{center}
\begin{tabular}{rrrrrr}
\toprule
$a$&$b$&$z$&$C_k(z)$&$u_a$&$u_b$\\
\midrule
40&41&56&91&184&193\\
40&41&57&90&184&193\\
41&40&56&91&193&184\\
41&40&57&90&193&184\\
\bottomrule
\end{tabular}
\end{center}
They are covered by a single coarse family H1, $a=40,41$, with
\begin{gather}
 \delta\le1,\quad R_p+R_q+\delta\le450,\notag\\
 \varepsilon_p\ge240+9(a-40),\quad
 \varepsilon_q\ge249-9(a-40),\notag\\
 S_p+S_q+\delta-2K\le1842,\qquad K=|2D|.\label{eq:H}
\end{gather}
These are consequences of \eqref{eq:joint}, $z\le57$, $w\le228$,
and the displayed lower mass bounds. The model does not retain a $z$
variable in H1. Since $\delta\le1<\min(a,b)+5$, ordinary quantiles apply.

\section{The complete common model}
\subsection{Modes and all parameter ranges}
The twenty center labels are inherited from \cite{base99,base98}; for
completeness define the outer magnitudes
$p_1>p_2>p_3>p_4>M_p$ and their $q$ reflection, and
\[
 \lambda_p=p_1-p_2-p_3-p_4,\quad
 \alpha_p=p_2-p_3-p_4,\quad\beta_p=2p_2-p_1-p_3-p_4.
\]
The zero-based positive centers are
\[
 v^p=(\lambda_p+p_1,\ldots,\lambda_p+p_4,
       \lambda_p-q_1,\ldots,\lambda_p-q_4,\alpha_p,\beta_p).
\]
Interchange $p,q$ for negative centers. Strings \texttt{p0},\ldots,
\texttt{q9} refer to these slots, not the magnitude variables.
Pair 1 uses \texttt{p4,q4}. Branch A means $\beta_p+\beta_q\le0$ and
pair 2 uses \texttt{p9,q9}; branch B means $\beta_p+\beta_q>0$ and
pair 2 uses \texttt{p5,q5}. The other sixteen units are singles.

Choose every $j$-element subset of missing units. Retain the other
singles; if pair 1 is present reflect to choose its positive contributor.
If pair 2 is present and cannot split, enumerate both contributor choices.
It may split precisely in the relaxed branch-A handling where a single
\texttt{p5} or \texttt{q5} is missing. In that case drop its individual
loss rows, choose an \emph{actually contributing} member $s=p$ or $q$,
and retain the necessary row
\begin{equation}\label{eq:member}
 -\beta_s-M_s\le-1.
\end{equation}
If pair 2 itself is missing there is no such choice. The foundation's
pairing lemma proves exhaustiveness: pair 1 and branch-B pair 2 have at
most one contributor; in branch A, two active singles \texttt{p5,q5}
force the same property for pair 2. Dropping a possibly split pair
weakens the loss sum, and choosing an actual contributor validates
\eqref{eq:member}. This row alone is not claimed to imply contribution.
The mode counts for $j=0,1,2$ are 4, 70, and 578.

The root list is completely specified by the following table. Each
large-loss family designates each retained center once. The 70 one-empty
modes have 1,186 retained centers in total. Every integer $a$ is included;
$b=81-a$ throughout.
\begin{center}\small
\begin{tabular}{lrrlr}
\toprule
Family&$j$&Range of $a$&Quantiles&Models\\
\midrule
Z0&0&26--55&Global&120\\
P0&0&27--54&Global&112\\
Z1&1&29--52&Core&1680\\
P1&1&30--51&Core&1540\\
ZL&1&34--47&One large center global; others core&16604\\
PL&1&36--45&One large center global; others core&11860\\
Z2&2&34--47&Core&8092\\
P2&2&36--45&Ordinary&5780\\
H1&1&40--41&Ordinary&140\\
\midrule
Total&&&&45928\\
\bottomrule
\end{tabular}
\end{center}
Z uses \eqref{eq:Z}, P uses \eqref{eq:P}, H uses \eqref{eq:H}.
The low-hole ranges follow from nonnegative loss budgets, and the L ranges
also impose $\delta\ge86$. The scalar checker independently verifies
that each of the 416 surviving low-hole states is covered by the table.

\subsection{Orders, mass, span, and loss rows}
All variables are nonnegative real numbers. The base variables are
\[
 p_1,\ldots,p_4,q_1,\ldots,q_4,M_p,M_q,R_p,R_q,S_p,S_q,\delta,
 \quad(l_v,u_v)_{v\text{ retained}}.
\]
Every family also exposes the four inner extremes:
\[
 x_0=R_p<x_1<x_2<x_3<x_4=M_p,
 \qquad y_0=R_q<y_1<y_2<y_3<y_4=M_q.
\]
They have integral gaps at least one. Include the positive-side rows
below and their reflections:
\begin{align*}
 p_{i+1}-p_i&\le-1\quad(1\le i\le3),&M_p-p_4&\le-1,\\
 R_p-M_p&\le-4,&-R_p&\le-a,\\
 x_{i-1}-x_i&\le-1\quad(1\le i\le4).&&
\end{align*}
The eight outer right sides become zero when $j=2$.
With zero empty units, all four own centers on each side contribute and
are integral, and
$\lambda_p=-v^p_0+v^p_1+v^p_2+v^p_3$ makes the outer points integral.
With one empty unit use the inherited contributing-center integrality
lemma. With two empty units we assume only weak outer order.
The branch row is $\beta_p+\beta_q\le0$ in A;
in B use $-\beta_p-\beta_q\le-1$ for $j\le1$ and $\le0$ for $j=2$.
When a split member is chosen, add \eqref{eq:member}.

Use all four mass rows, separately, on each side:
\begin{align}
 S_p&\ge T_a+\varepsilon_p+3R_p+6+M_p,\notag\\
 S_p&\ge T_{a-1}+4R_p+6+M_p,\notag\\
 S_p&\ge T_a+\varepsilon_p+\sum_{i=1}^4x_i,\notag\\
 S_p&\ge T_{a-1}+R_p+\sum_{i=1}^4x_i.\label{eq:mass}
\end{align}
They follow respectively from core mass and three inner gaps, distinct
core integers and those gaps, core mass with exposed inner points, and
distinct core integers with exposed points. Replace $\varepsilon_p$
by its affine lower bound in the family; it is not a new variable.
The first row is multiplied by two in the implementation.
These lower bounds are imposed alongside each other, not added together.

Include the family's loss and joint-span budgets, \eqref{eq:loss}, the
quantile rows in the table, and the slot caps from the loss classification.
In an L model add $-l_v-u_v\le-86$ for its designated center.
In P models include $-z\le-1$, $z\le14$ and the same $z$ in every
mass, loss, span, and cardinality row.

\subsection{Five nested envelopes and one coarse sumset bound}
Set $K=|2D|$, $W_t=x_t+y_t$, $s_t=2(4-t)$.
The nested-envelope argument of \cite{base98} gives for $0\le t\le4$
\begin{equation}\label{eq:envelope}
 K\le2W_t+1+\binom{s_t+1}{2}
 +\sum_{i=t+1}^4\!\left[
 \min(W_t+1,x_i-x_{i-1})+\min(W_t+1,y_i-y_{i-1})\right].
\end{equation}
For clarity, replace the inner set through level $t$ by its integer hull
$[-y_t,x_t]$. Its double sum has $2W_t+1$ points. Adding each successive
external translate enlarges the union by at most the minimum of the
hull length $W_t+1$ and the new gap. Only pairs of the $s_t$ external
points remain, at most $\binom{s_t+1}{2}$ with repetition. This proves
\eqref{eq:envelope}, including $t=4$, where the sum is empty and
$K\le2M_p+2M_q+1$.

For $t<4$ introduce nonnegative variables $g^p_{ti},g^q_{ti}$ bounded
above by the two arguments of the corresponding minimum:
\[
 g^p_{ti}-x_i+x_{i-1}\le0,\qquad g^p_{ti}-x_t-y_t\le1,
 \qquad t<i\le4,
\]
and the reflected rows. In \eqref{eq:envelope} replace each minimum by
its $g$ variable. An actual set realizes the model by choosing the minima;
we need not force equality with the minima in the relaxation.
Also retain the coarse bound
\begin{equation}\label{eq:coarse}
 K-M_p-M_q-R_p-R_q\le21.
\end{equation}
All pairs except two positive outliers or two negative outliers lie in
$[-M_q-R_q,M_p+R_p]$; there are at most ten repeated pairs of each
same-sign outlier type. The hull has $M_p+M_q+R_p+R_q+1$ positions,
proving \eqref{eq:coarse}.
Finally, the definition of $c$ gives the exact identity
\[
 S_p+S_q+\delta-2K=18d+84-jd+w.
\]
Thus the cardinality row is
\begin{equation}\label{eq:cardinality}
 S_p+S_q+\delta-2K\le1704-90j+4z
\end{equation}
for low holes (set $z=0$ in Z), and is the row in \eqref{eq:H} for H1.
These are all rows of the delivered model. The slot caps, $t=4$ envelope,
and coarse bound are valid strengthenings, not solver-inferred cuts.

\section{Exact certification, checks, and the next endpoint}
For each structural group the matrix $A$ is constant. Its right-hand
side is a rational quadratic polynomial in $a$, because
$b=81-a$, $I=(a-40)(a-41)$, and $T_a,T_b$ are quadratic.
The shared hole parameter in P families is a model coordinate, not an
interpolated parameter. For a nonnegative integer row combination $w$,
verify $wA\ge0$ coordinatewise and maximize $p(a)=wb(a)$ over the assigned
integer interval. Endpoints suffice for a convex or linear polynomial;
for a concave quadratic also check the floor and ceiling of its vertex
when they lie in the interval. Strict negativity proves inconsistency
for every integer $a$ in that interval.
The polynomial class builds the formulas symbolically with rational
coefficients and rejects degrees above two. It does not infer formulas
from sample values.

The certificate pool contains 209 integer vectors. Complete coverage has
3,746 structural groups and 3,750 assigned intervals, covering all 45,928
models. The checker independently reconstructs the mode and parameter
universe and rejects missing groups, gaps, overlaps, out-of-range
assignments, or unused vectors. The least normalized contradiction
$-\max p(a)/\sum_iw_i$ is $12/349$.
The 3,750 matrix-combination checks replace 45,928 repetitions, a reduction
of about 91.8\%; comparison across different endpoint universes should
not be interpreted as an identical-workload speed benchmark.

The \texttt{--exhaustive} mode was also run: it rebuilds every numerical
model, compares every coefficient and right-hand side with the symbolic
formulas, verifies each integer certificate, and compares its minimum
margin with the exact interval extremum. Both results agree apart from
the mode label. The regenerated certificate file is byte-identical.
Standard and \texttt{-O} outputs are byte-identical;
the previous upper-98 finite result is reproduced unchanged. Verification
uses explicit checks, not Python assertions that disappear under
optimization. Ten negative controls reject bad weights, noncontradictory
weights, unknown rows, an interior-vertex failure, missing coverage,
a missing actual-contributor choice, an invalid flexible threshold,
the failed coarse lattice gate, and an attempt to certify size 97.
A further 5,409 exact inverse-boundary checks corroborate
\eqref{eq:flex}; the general proof is Lemma~\ref{lem:flex}.

\begin{proof}[Proof of Theorem~\ref{thm:main}]
The foundation excludes all sizes above 98. A hypothetical size-98 set
passes the sign bootstrap, inverse gates, and refined lattice argument.
High-hole all-active configurations are excluded by the mass comparisons;
all high-hole states with an empty unit map into H1. Low holes map into
the eight complete loss families by the integer loss threshold and the
number of empty units. Every actual set satisfies all rows of one model,
including the shared hole parameter, actual contributor choice, and
valid strengthenings. Every such model has a strict exact integer-row
contradiction. Hence no size-98 feasible set exists.
\end{proof}

\clearpage
The delivered files are:
\begin{description}
\item[Checker] \path{verification/upper97.py}.
\item[Proof data] \path{data/upper97_certificates.json}.
\item[Generator] \path{experiments/generate_upper97_certificates.py}.
\item[Report] \path{reports/upper97_from98_20261016.tex} and its PDF.
\item[Run records] \path{results/2026-10-16/upper97/}.
\end{description}
The records contain full finite results, input hashes, baseline output,
and a run receipt.
From the repository root, with Python 3.10 or newer:
\begin{verbatim}
python -I -S -B verification/upper97.py --output .work/upper97.json
python -I -S -B verification/upper97.py --exhaustive
python -I -S -B verification/upper98.py
\end{verbatim}
No third-party package, optimizer, or network is needed for checking.
The optional generator uses SciPy to propose weights and accepts them
only after exact reconstruction and verification. The Windows entry
\texttt{verify\_upper97\_windows.bat} runs the five continuation stages.
The printed proof supplies the infinite structural arguments; an
arithmetic PASS is not a proof-assistant formalization of those arguments.

The next endpoint is size 97. Its balanced zero-hole budget permits
$\delta\le70$ after two empty units. At positive low holes it permits
$\delta\le46$, which can exceed the minimum outside loss 45.
Thus the ordinary quantiles used in P2 here cannot simply be reused.
No exclusion of size 97 is claimed by this certificate set, and the
stronger archived bound 84 remains subject to its separate complete audit.

\begin{thebibliography}{9}
\bibitem{base103} Sum-of-Subset project,
\emph{A checked continuation from 115 to 103}, 14 October 2026.
\path{reports/upper103_from115_20261014.tex}.
\bibitem{base100} Sum-of-Subset project,
\emph{A checked continuation from 103 to 100}, 16 October 2026.
\path{reports/upper100_from103_20261016.tex}.
\bibitem{base99} Sum-of-Subset project,
\emph{A checked continuation from 100 to 99}, 16 October 2026.
\path{reports/upper99_from100_20261016.tex}.
\bibitem{base98} Sum-of-Subset project,
\emph{An optimized proof path from 99 to 98}, 16 October 2026.
\path{reports/upper98_from99_20261016.tex}.
\bibitem{archive} Sum-of-Subset project,
\emph{Complete upper-bound 86 report}, 5 October 2026,
section ``Two further exclusions: the general upper bound 96'' and its
certificate appendix. \path{reports/upper86_complete_20261005.tex}.
\end{thebibliography}
\end{document}
